\section{Method}
\label{sec:method}

\subsection{The \textsc{IRT-Fit} Framework}
\label{sec:framework}

Our method takes as input a binary response matrix $\mathbf{X} \in \{0,1\}^{M \times n}$, where $M$ models answer $n$ benchmark items, together with pre-calibrated four-parameter logistic (4PL) item parameters from PSN-IRT \citep{psnirt2026}.
The 4PL model defines the probability that a model with ability $\theta$ answers item $i$ correctly:
\begin{equation}
  P_i(\theta) = c_i + (d_i - c_i)\,\frac{1}{1 + e^{-a_i(\theta - b_i)}}
  \label{eq:4pl}
\end{equation}
where $a_i$ is discrimination, $b_i$ is difficulty, $c_i$ is a lower asymptote (guessing), and $d_i$ is an upper asymptote (slip).

The framework proceeds in three steps.
\textbf{Step~1}: Estimate each model's ability $\hat{\theta}$ via maximum likelihood on each benchmark separately.
\textbf{Step~2}: Compute the person-fit statistic $\ell_z^*$ \citep{snijders2001}, a standardized likelihood residual that quantifies how well a model's item-level response pattern conforms to the IRT model given $\hat{\theta}$.
The statistic extends the original $\ell_z$ \citep{drasgow1985} by correcting for estimation uncertainty in $\hat{\theta}$, so that $\ell_z^* \sim \mathcal{N}(0,1)$ under the null hypothesis.
\textbf{Step~3}: For within-model contamination detection, compute $\Delta\ell_z^* = \ell_z^*(\text{target}) - \ell_z^*(\text{reference})$ between a target benchmark and a reference benchmark from the same model.
The entire procedure is unsupervised: it requires no training data, no labeled contamination examples, and no access to pretraining corpora.

\paragraph{Validity under assumption violations.}
IRT assumes local independence and unidimensionality, conditions that LLM benchmarks violate.
When these assumptions fail, $\ell_z^*$ can exhibit systematic bias: all models may appear aberrant even without contamination, because inter-item dependencies inflate the variance of response patterns beyond what the IRT model predicts.
However, this bias is approximately constant across benchmarks for a given model, because the same latent capabilities and the same dependency structure drive responses.
The $\Delta\ell_z^*$ statistic functions as a difference-in-differences (DiD) estimator: the shared bias cancels in the subtraction, isolating the contamination-specific signal.
We calibrate significance thresholds via a nonparametric bootstrap over the null distribution of $\Delta\ell_z^*$ rather than relying on the asymptotic $\mathcal{N}(0,1)$ reference, which accommodates residual model misspecification.
Universal misfit (all models showing $|\ell_z^*| > 1.96$) is therefore a finding about benchmark properties, not an artifact; $\Delta\ell_z^*$ bypasses absolute thresholds entirely.

\subsection{Detection Power Characterization}
\label{sec:power_theory}

We derive the statistical power of $\ell_z^*$ as a function of benchmark size $n$ and contamination fraction $\pi$.

\paragraph{Contamination model.}
A fraction $\pi$ of items are forced correct regardless of ability.
The expected response becomes:
\begin{equation}
  \mathbb{E}[x_i \mid \theta_0, \pi] = P_i(\theta_0) + \pi \cdot Q_i(\theta_0)
  \label{eq:contam_response}
\end{equation}
where $Q_i(\theta) = 1 - P_i(\theta)$.
The contaminated response inflates the MLE: $\mathbb{E}[\hat{\theta}] > \theta_0$ for $\pi > 0$, because the score equation $\sum_i (\mathbb{E}[x_i] - P_i(\theta))\,h_i(\theta) = 0$ shifts the solution upward.

\paragraph{Proposition 1 ($\sqrt{n}$ scaling).}
The noncentrality parameter $\delta(\theta_0, \pi)$ of $\ell_z^*$ under contamination satisfies $\delta \propto \sqrt{n}$.
The numerator accumulates $n$ per-item residuals ($O(n)$) while the denominator grows as $\sqrt{\mathrm{Var}^*} = O(\sqrt{n})$ from the \citet{snijders2001} variance formula, so
\begin{equation}
  \delta(\theta_0, \pi) = \sqrt{n}\;\cdot\;\delta_1(\theta_0, \pi)
  \label{eq:sqrt_n}
\end{equation}
where $\delta_1$ is the per-item noncentrality that depends only on item parameters and $(\theta_0, \pi)$.

\paragraph{Proposition 2 (marginal power).}
Integrating over $\theta_0 \sim \mathcal{N}(0,1)$ via Gauss--Hermite quadrature ($K{=}30$ nodes) yields the marginal power under two-sided testing at level $\alpha$:
\begin{equation}
  \mathrm{Power}(n, \pi) = \sum_{k=1}^{K} w_k \left[\Phi\!\left(\frac{-z_\alpha - \delta_k}{\sigma_k}\right) + 1 - \Phi\!\left(\frac{z_\alpha - \delta_k}{\sigma_k}\right)\right]
  \label{eq:power}
\end{equation}
where $\delta_k$ and $\sigma_k$ are the conditional mean shift and standard deviation of $\ell_z^*$ at the $k$-th quadrature node, $w_k$ are the quadrature weights, and $z_\alpha = 1.96$.
Two-sided testing is essential: at moderate contamination ($\pi < 80\%$), $\ell_z^*$ is negative (aberrant pattern), but at high contamination ($\pi > 85\%$), $\ell_z^*$ turns positive (over-consistent pattern) as the inflated $\hat{\theta}$ renders near-perfect accuracy expected.
The crossover occurs at $\pi^* \approx 81$--$85\%$ across benchmarks.
Practitioners can determine minimum benchmark sizes from Equation~\ref{eq:power}: for example, 80\% power at 15\% contamination requires $n \geq 7{,}958$ items (derived under MMLU item parameters; the threshold varies across benchmarks due to differences in per-item discriminatory power).

\subsection{$\theta$-Shift Analysis}
\label{sec:theta_shift}

Contamination creates an identification problem: higher accuracy raises $\hat{\theta}_{\text{MLE}}$, which in turn ``explains away'' the aberrant responses and attenuates the $\ell_z^*$ signal.
We term this the \emph{$\theta$-shift mechanism}.
As $\pi$ increases, $\hat{\theta}$ absorbs a larger share of the contamination signal into a legitimate ability increase, so the free-$\theta$ dose-response curve becomes non-monotonic: $|\Delta\ell_z^*|$ can decrease at higher contamination levels.

The solution is \emph{$\theta$-fixed analysis}: anchor ability at a pre-contamination baseline $\theta_{\text{fixed}}$ estimated from clean conditions or cross-validated from held-out benchmarks.
With $\theta$ fixed, the $\ell_z^*$ statistic captures the full residual between observed and expected performance, and the dose-response curve becomes monotonically increasing.
In our controlled experiments (\S\ref{sec:ctrl_contam}), fixed-$\theta$ estimation yields a monotonic signal ratio of 3.12 between 50\% and 15\% contamination, whereas free-$\theta$ estimation produces an inverted ratio of 0.29.

For practical deployment, $\theta_{\text{fixed}}$ can be obtained through two strategies: (1) cross-benchmark estimation, where $\theta$ is estimated from benchmarks presumed clean and applied to the target benchmark; or (2) temporal anchoring, where $\theta$ from an earlier evaluation serves as the reference for subsequent evaluations of the same model.